\section{Local Gauss operators of a finite group}
\label{sec:monomial_gauss}

\begin{lemma}[Finite-group unitary average is an orthogonal projection]
    \label{thm:mpug_finite_group_unitary_average_projection}
    \lean{Matrix.isStarProjection_inv_card_smul_sum,
        Matrix.isOrthogonalProjection_inv_card_smul_sum}
    \leanok
    Let $G$ be a finite group and let $\rho\colon G\to\mathrm U(n)$ be a
    unitary representation. The normalized average
    \begin{align}
        P_\rho & = \frac{1}{|G|}\sum_{g\in G}\rho(g)
        \label{eq:mpug_finite_group_unitary_average}
    \end{align}
    satisfies
    \begin{align}
        P_\rho^\dagger & =P_\rho, \\
        P_\rho^2       & =P_\rho.
    \end{align}
    Hence $P_\rho$ is an orthogonal projection. For the local Gauss operators
    of \cite[lines~457--461]{FrancoRubio2025MPUGauging} it is the local
    Gauss-law average.
\end{lemma}

\begin{proof}\leanok
    When each matrix $\rho(g)$ is viewed as the corresponding linear map,
    $P_\rho$ is the averaging operator of the representation, which is
    idempotent. For self-adjointness, $\rho(g)^\dagger=\rho(g^{-1})$, and
    reindexing the finite sum by $g\mapsto g^{-1}$ gives
    $P_\rho^\dagger=P_\rho$.
\end{proof}

\begin{definition}[Unitary factorization comparison]
    \label{def:mpug_unitary_factorization_comparison}
    \lean{Matrix.unitaryFactorizationComparison,
        Matrix.unitaryFactorizationComparison_coe}
    \leanok
    Let $I$ be a set and let $R_{a,b}$ be a family of unitary matrices of a
    fixed size, indexed by $(a,b)\in I^2$. Following the orientation in the
    corollary after Proposition~3
    of~\cite[lines~2782--2821]{FrancoRubio2025MPUGauging}, the unitary
    comparison from $(a,b)$ to the reference pair $(c,d)$ is
    \begin{align}
        \lambda_{a,b}^{c,d}
         & =(R_{c,d})^{-1}R_{a,b}
        =R_{c,d}^\dagger R_{a,b}.
        \label{eq:mpug_unitary_factorization_comparison}
    \end{align}
    In that corollary, $I=G$ and comparisons are used when $ab=cd$; this
    equality selects the factorizations and is not needed in the unitary
    calculation.
\end{definition}

\begin{lemma}[Identity, inverse, composition, and right-reference laws]
    \label{thm:mpug_unitary_factorization_comparison_laws}
    \lean{Matrix.unitaryFactorizationComparison_self,
        Matrix.unitaryFactorizationComparison_inv,
        Matrix.unitaryFactorizationComparison_trans,
        Matrix.unitaryFactorizationComparison_of_reference_eq_one,
        Matrix.unitaryFactorizationComparison_right_reference}
    \leanok
    \uses{def:mpug_unitary_factorization_comparison}
    For all indices $a,b,c,d,f,g\in I$,
    \begin{align}
        \lambda_{a,b}^{a,b} & =\mathbf 1.
    \end{align}
    Reversing the comparison gives
    \begin{align}
        \left(\lambda_{a,b}^{c,d}\right)^{-1}
         & =\lambda_{c,d}^{a,b}.
    \end{align}
    Comparisons compose as
    \begin{align}
        \lambda_{a,b}^{f,g}
         & =\lambda_{c,d}^{f,g}\lambda_{a,b}^{c,d}.
        \label{eq:mpug_unitary_factorization_comparison_laws}
    \end{align}
    If $R_{c,d}=\mathbf 1$, then
    $\lambda_{a,b}^{c,d}=R_{a,b}$. In particular, if $I=G$ is a group, the
    normalization $R_{e,x}=\mathbf 1$ gives the right-reference specialization
    $\lambda^R_{g,h}=\lambda_{g,h}^{e,gh}$ in the cited corollary,
    \begin{align}
        R_{g,h} & =\lambda_{g,h}^{e,gh}.
        \label{eq:mpug_unitary_factorization_right_reference}
    \end{align}
\end{lemma}

\begin{proof}\leanok
    \uses{def:mpug_unitary_factorization_comparison}
    The identity, inverse, and right-reference formulas follow from
    $R_{a,b}^{-1}R_{a,b}=\mathbf 1$, inversion of a product, and
    $\mathbf 1^{-1}R_{a,b}=R_{a,b}$. For composition, inserting the identity at
    the intermediate pair gives
    \begin{align}
        (R_{f,g})^{-1}R_{a,b}
         & =(R_{f,g})^{-1}R_{c,d}(R_{c,d})^{-1}R_{a,b}.
    \end{align}
\end{proof}

\begin{definition}[Gauss action on pairs of gauge labels]
    \label{def:mpug_gauss_leg_action}
    \lean{GaussLaw.gaussLegAction,
        GaussLaw.gaussLegAction_apply}
    \leanok
    Let $G$ be a group. Following the two gauge labels in the local Gauss
    operator of \cite[lines~3380--3387]{FrancoRubio2025MPUGauging}, for
    $g\in G$ define the permutation $T_g\colon G^2\to G^2$ by
    \begin{align}
        T_g(a,b) & =\left(ag^{-1},gb\right).
        \label{eq:mpug_gauss_leg_action}
    \end{align}
\end{definition}

\begin{lemma}[Composition of the Gauss action]
    \label{thm:mpug_gauss_leg_action_mul}
    \lean{GaussLaw.gaussLegAction_one,
        GaussLaw.gaussLegAction_mul}
    \leanok
    \uses{def:mpug_gauss_leg_action}
    For all $g,h\in G$,
    \begin{align}
        T_e          & =\id,    \\
        T_g\circ T_h & =T_{gh}.
        \label{eq:mpug_gauss_leg_action_mul}
    \end{align}
\end{lemma}

\begin{proof}\leanok
    \uses{def:mpug_gauss_leg_action}
    The two coordinates give
    \begin{align}
        T_g(T_h(a,b))
         & =\left(ah^{-1}g^{-1},ghb\right)
        =\left(a(gh)^{-1},(gh)b\right).
    \end{align}
\end{proof}

\begin{definition}[Local Gauss operator]
    \label{def:mpug_gauss_operator}
    \lean{GaussLaw.gaussOperator,
        GaussLaw.gaussOperator_apply_target,
        GaussLaw.gaussOperator_apply_of_ne}
    \leanok
    \uses{def:mpug_unitary_factorization_comparison,
        def:mpug_gauss_leg_action}
    Let $G$ be finite, let $R_{a,b}\in\mathrm U(n)$ for $a,b\in G$, and take
    the basis $\ket{i,a,b}$ of $\C^n\otimes\C[G]^{\otimes2}$. As in
    \cite[lines~3380--3387]{FrancoRubio2025MPUGauging}, define the local Gauss
    operator
    \begin{align}
        \mathcal G_g
         & =\sum_{a,b\in G}
        \lambda_{a,b}^{ag^{-1},gb}
        \otimes\ket{ag^{-1},gb}\!\bra{a,b}.
        \label{eq:mpug_gauss_operator}
    \end{align}
    Equivalently, the block from $(a,b)$ to $T_g(a,b)$ is
    $(R_{T_g(a,b)})^{-1}R_{a,b}$ and every other block in that column
    vanishes.
\end{definition}

\begin{lemma}[Gauss operators form a unitary representation]
    \label{thm:mpug_gauss_representation}
    \lean{GaussLaw.gaussOperator_one,
        GaussLaw.gaussOperator_mul,
        GaussLaw.star_gaussOperator,
        GaussLaw.gaussOperator_mem_unitaryGroup,
        GaussLaw.gaussRepresentation,
        GaussLaw.gaussRepresentation_coe}
    \leanok
    \uses{def:mpug_gauss_operator}
    The local Gauss operators satisfy
    \begin{align}
        \mathcal G_e & =\id.
    \end{align}
    Their product law is
    \begin{align}
        \mathcal G_g\mathcal G_h & =\mathcal G_{gh}.
        \label{eq:mpug_gauss_representation}
    \end{align}
    Finally,
    \begin{align}
        \mathcal G_g^\dagger & =\mathcal G_{g^{-1}}.
    \end{align}
    Hence $g\mapsto\mathcal G_g$ is a unitary representation, establishing the
    representation part of the local Gauss-law proposition in
    \cite[lines~3380--3387]{FrancoRubio2025MPUGauging}; the invariance
    statement of that proposition is not included here.
\end{lemma}

\begin{proof}\leanok
    \uses{thm:mpug_unitary_factorization_comparison_laws,
        thm:mpug_gauss_leg_action_mul,
        def:mpug_gauss_operator}
    In a column labelled by $(a,b)$, multiplication has only the intermediate
    gauge label $T_h(a,b)$. The surviving block is
    \begin{align}
        \lambda_{T_h(a,b)}^{T_g(T_h(a,b))}
        \lambda_{a,b}^{T_h(a,b)}
         & =\lambda_{a,b}^{T_{gh}(a,b)},
    \end{align}
    by the comparison composition law and
    $T_gT_h=T_{gh}$. At $e$, every surviving comparison is the identity.
    Reversing a comparison gives
    $(\lambda_{a,b}^{c,d})^\dagger=\lambda_{c,d}^{a,b}$, so taking the adjoint
    gives $\mathcal G_g^\dagger=\mathcal G_{g^{-1}}$. The product with this
    adjoint is $\mathcal G_e=\id$, which proves unitarity.
\end{proof}

\begin{definition}[Local Gauss projector]
    \label{def:mpug_gauss_projector}
    \lean{GaussLaw.gaussProjector,
        GaussLaw.gaussProjector_eq_average}
    \leanok
    \uses{thm:mpug_gauss_representation}
    Let $G$ be finite and let $\mathcal G_g$ act on
    $\mathcal H=\C^n\otimes\C[G]^{\otimes2}$. Define the local projector
    \begin{align}
        P_G(R) & =\frac{1}{|G|}\sum_{g\in G}\mathcal G_g.
        \label{eq:mpug_gauss_projector}
    \end{align}
    This is the local average introduced in
    \cite[lines~3494--3499]{FrancoRubio2025MPUGauging}, before it is placed
    at a position of a chain.
\end{definition}

\begin{lemma}[The local Gauss average is a projection]
    \label{thm:mpug_gauss_projector_projection}
    \lean{GaussLaw.isStarProjection_gaussProjector}
    \leanok
    \uses{def:mpug_gauss_projector}
    The operator in~(\ref{eq:mpug_gauss_projector}) satisfies
    \begin{align}
        P_G(R)^2       & =P_G(R), \\
        P_G(R)^\dagger & =P_G(R).
    \end{align}
\end{lemma}

\begin{proof}\leanok
    \uses{thm:mpug_finite_group_unitary_average_projection,
        thm:mpug_gauss_representation}
    The operator $P_G(R)$ is the normalized average of the unitary
    representation $g\mapsto\mathcal G_g$, which is an orthogonal projection
    by Lemma~\ref{thm:mpug_finite_group_unitary_average_projection}.
\end{proof}

\begin{lemma}[Local Gauss operator of monomial unitaries]
    \label{thm:mpug_gauss_operator_monomial}
    \lean{GaussLaw.gaussMonomialPerm,
        GaussLaw.gaussMonomialPerm_apply,
        GaussLaw.gaussMonomialPhase,
        GaussLaw.gaussOperator_eq_monomial}
    \leanok
    \uses{def:mpug_gauss_operator, def:mpug_monomial_matrix}
    Suppose that every unitary of the family $R$ is monomial,
    $R_{a,b}=M(\sigma_{a,b},\varphi_{a,b})$. Then the local Gauss operator
    of Definition~\ref{def:mpug_gauss_operator} is the monomial operator
    \begin{align}
        \mathcal G_g\ket{i,(a,b)}
         & =\overline{\varphi_{T_g(a,b)}\bigl(\sigma_{T_g(a,b)}^{-1}\sigma_{a,b}(i)\bigr)}\,
        \varphi_{a,b}(i)\,
        \ket{\sigma_{T_g(a,b)}^{-1}\sigma_{a,b}(i),\ T_g(a,b)}.
        \label{eq:mpug_gauss_operator_monomial}
    \end{align}
\end{lemma}

\begin{proof}\leanok
    \uses{thm:mpug_monomial_matrix_laws, def:mpug_gauss_operator}
    The block of $\mathcal G_g$ from $(a,b)$ to $T_g(a,b)$ is
    $R_{T_g(a,b)}^\dagger R_{a,b}$, which is monomial by
    (\ref{eq:mpug_monomial_mul}) and (\ref{eq:mpug_monomial_adjoint}); all
    other blocks vanish.
\end{proof}

\begin{definition}[Prescribed defect maps]
    \label{def:mpug_defect_maps}
    \lean{GaussLaw.DefectMaps}
    \leanok
    Let $\mathcal H_{\mathrm m}$ be a finite-dimensional matter space.
    Prescribed defect maps assign to every ordered pair $(a,b)\in G^2$ a defect
    subspace
    \begin{align}
        \mathcal K_{a,b} & \subseteq\mathcal H_{\mathrm m}
    \end{align}
    together with a map $D_{a,b}$ of $\mathcal H_{\mathrm m}$ whose
    restriction to that subspace is an isometry: for every
    $\ket{\xi}\in\mathcal K_{a,b}$,
    \begin{align}
        D_{a,b}^\dagger D_{a,b}\ket{\xi}
         & =\ket{\xi}.
        \label{eq:mpug_defect_isometry}
    \end{align}
    Such maps are prescribed only on these subspaces in the construction
    of~\cite[lines~4215--4326]{FrancoRubio2025MPUGauging}.
\end{definition}

\begin{definition}[Full unitary completion class]
    \label{def:mpug_completion_class}
    \lean{GaussLaw.DefectMaps.IsCompletion,
        GaussLaw.DefectMaps.completionClass,
        GaussLaw.DefectMaps.mem_completionClass_iff}
    \leanok
    \uses{def:mpug_defect_maps}
    The full completion class of prescribed defect maps
    $D=(D_{a,b})_{a,b\in G}$ is
    \begin{align}
        \mathfrak C(D)
         & =\left\{R=(U_{a,b})_{a,b\in G}:
        U_{a,b}\in\mathrm U(\mathcal H_{\mathrm m}),\
        U_{a,b}|_{\mathcal K_{a,b}}=D_{a,b}|_{\mathcal K_{a,b}}\right\}.
        \label{eq:mpug_completion_class}
    \end{align}
    As noted in~\cite[lines~4215--4326]{FrancoRubio2025MPUGauging}, when
    $\mathcal K_{a,b}$ is a proper subspace of the matter space, the unitary
    extension of $D_{a,b}$ is non-unique.
\end{definition}

\begin{lemma}[Existence of full unitary completions]
    \label{thm:mpug_completion_class_nonempty}
    \lean{GaussLaw.DefectMaps.completionClass_nonempty,
        GaussLaw.DefectMaps.exists_completion}
    \leanok
    \uses{def:mpug_defect_maps, def:mpug_completion_class}
    Every family of prescribed defect maps satisfying
    (\ref{eq:mpug_defect_isometry}) admits a simultaneous unitary extension:
    \begin{align}
        \mathfrak C(D) & \neq\varnothing.
        \label{eq:mpug_completion_class_nonempty}
    \end{align}
    Equivalently, there is a family
    $R=(U_{a,b})_{(a,b)\in G^2}$ of unitary operators such that, for every
    $a,b\in G$ and every $\ket{\xi}\in\mathcal K_{a,b}$,
    \begin{align}
        U_{a,b}\ket{\xi}
         & =D_{a,b}\ket{\xi}.
        \label{eq:mpug_completion_exists}
    \end{align}
    This is the finite-dimensional extension step implicit in
    \cite[lines~4215--4326]{FrancoRubio2025MPUGauging}, where the subspaces
    $\mathcal K_{a,b}$ and their prescribed isometries enter as hypotheses.
\end{lemma}

\begin{proof}\leanok
    \uses{def:mpug_defect_maps, def:mpug_completion_class}
    For each pair $(a,b)$, the restriction of $D_{a,b}$ to
    $\mathcal K_{a,b}$ is a linear isometry and therefore extends to a unitary
    operator on the whole finite-dimensional matter space. Choosing one such
    extension for every pair gives a family $R$ satisfying
    (\ref{eq:mpug_completion_exists}), hence a member of $\mathfrak C(D)$.
\end{proof}

\begin{lemma}[Comparison identity for a completion]
    \label{thm:mpug_completion_transport}
    \lean{GaussLaw.DefectMaps.IsCompletion.mulVec_eq_prescribed,
        GaussLaw.DefectMaps.IsCompletion.mulVec_eq_mulVec,
        GaussLaw.DefectMaps.IsCompletion.unitaryFactorizationComparison_mulVec,
        GaussLaw.DefectMaps.IsCompletion.unitaryFactorizationComparison_gaussLegAction_mulVec}
    \leanok
    \uses{def:mpug_completion_class,
        def:mpug_unitary_factorization_comparison,
        def:mpug_gauss_leg_action}
    Let $R=(U_{a,b})\in\mathfrak C(D)$, let
    $\ket{\xi}\in\mathcal K_{a,b}$ and $\ket{\eta}\in\mathcal K_{c,d}$, and
    assume the exact covariance
    \begin{align}
        D_{a,b}\ket{\xi} & =D_{c,d}\ket{\eta}.
        \label{eq:mpug_prescribed_covariance}
    \end{align}
    Then, as required for the modified Gauss laws
    of~\cite[lines~4325--4335]{FrancoRubio2025MPUGauging} when $(c,d)=T_g(a,b)$,
    \begin{align}
        U_{a,b}\ket{\xi}             & =D_{a,b}\ket{\xi},
        \label{eq:mpug_completion_agrees}                  \\
        U_{a,b}\ket{\xi}             & =U_{c,d}\ket{\eta},
        \label{eq:mpug_completion_covariance}              \\
        \lambda_{a,b}^{c,d}\ket{\xi} & =\ket{\eta}.
        \label{eq:mpug_completion_transport}
    \end{align}
\end{lemma}
\begin{proof}\leanok
    \uses{def:mpug_completion_class,
        thm:mpug_unitary_factorization_comparison_laws}
    Identity~(\ref{eq:mpug_completion_agrees}) is the defining property of a
    completion on the defect subspace. Applying it at $(a,b)$ and at $(c,d)$
    and inserting~(\ref{eq:mpug_prescribed_covariance})
    gives~(\ref{eq:mpug_completion_covariance}), an equality of vectors in the
    whole matter space. Multiplying that equality by $U_{c,d}^\dagger$ and
    using $U_{c,d}^\dagger U_{c,d}=\Id$ gives
    \begin{align}
        \lambda_{a,b}^{c,d}\ket{\xi}
         & =U_{c,d}^\dagger U_{a,b}\ket{\xi}
        =U_{c,d}^\dagger U_{c,d}\ket{\eta}
        =\ket{\eta}.
    \end{align}
    The adjoint is applied only after the equality in the whole matter space
    is available; no step asserts that $U_{c,d}^\dagger$ restricts to
    $D_{c,d}^\dagger$ on a defect subspace.
\end{proof}
